% ======================================================================
% Prefix-Aware Tree Construction: Formal Analysis
% ======================================================================
% Self-contained section suitable for inclusion in a paper or appendix.
% Compile with any standard LaTeX setup (amsmath, amssymb, amsthm).
% ======================================================================

\newtheorem{definition}{Definition}
\newtheorem{proposition}{Proposition}
\newtheorem{corollary}{Corollary}
\newtheorem{remark}{Remark}

% ======================================================================
\section{Prefix-Aware Tree Construction}
\label{sec:prefix-aware-tree}
% ======================================================================

We formalize the tree construction problem for speculative decoding,
show that the natural objective is monotone submodular, and derive an
approximation guarantee for our greedy algorithm.

% ----------------------------------------------------------------------
\subsection{Setup and Notation}
\label{sec:setup}
% ----------------------------------------------------------------------

Consider a draft model that produces, at each depth $d = 1, \dots, D$,
a probability distribution $p_d(\cdot)$ over the vocabulary~$\mathcal{V}$.
A \emph{verification tree} $T$ is specified by a set of \emph{leaves},
where each leaf $\ell = (\sigma_1, \dots, \sigma_{D_\ell})$ is a token
sequence of length $D_\ell \leq D$.  The tree is packed into a trie and
verified in a single target-model forward pass.

During verification, the target model produces a deterministic output
sequence given the context.  The \emph{accepted length} $\tau(T)$ is
the length of the longest prefix of any leaf that matches the target's
output.  Because the target output is deterministic for a given context,
$\tau(T)$ is itself deterministic.

To construct the tree \emph{before} seeing the target output, we
optimize a surrogate objective defined by the draft probabilities.

% ----------------------------------------------------------------------
\subsection{The Prefix-Coverage Objective}
\label{sec:objective}
% ----------------------------------------------------------------------

\begin{definition}[Prefix-coverage objective]
\label{def:objective}
For a leaf $\ell = (\sigma_1, \dots, \sigma_{D_\ell})$, define its
\emph{depth-$k$ prefix} as $\pi_k(\ell) = (\sigma_1, \dots, \sigma_k)$
for $k \leq D_\ell$.  For a set of leaves $T$, let
\[
  \mathrm{Prefixes}_k(T)
  = \bigl\{\, \pi_k(\ell) : \ell \in T,\; D_\ell \geq k \,\bigr\}
\]
be the set of \emph{distinct} depth-$k$ prefixes covered by~$T$.
Define the \emph{prefix-coverage objective}
\begin{equation}
\label{eq:objective}
  f(T)
  \;=\; \sum_{k=1}^{D}\;\sum_{\sigma \,\in\, \mathrm{Prefixes}_k(T)}
        w(\sigma),
  \qquad
  w(\sigma)
  \;=\; \prod_{d=1}^{|\sigma|} p_d(\sigma_d).
\end{equation}
\end{definition}

\begin{remark}[Interpretation as expected accepted length]
\label{rem:interpretation}
The objective $f(T)$ equals the expected accepted length
$\mathbb{E}[\tau(T)]$ under a specific probabilistic model: at each
depth~$d$, the target token is drawn independently from the draft
distribution~$p_d$.  Under this \emph{draft-induced product model},
\[
  \mathbb{E}[\tau(T)]
  = \sum_{k=1}^{D} \Pr[\tau(T) \geq k]
  = \sum_{k=1}^{D}\;\sum_{\sigma \,\in\, \mathrm{Prefixes}_k(T)}
    \prod_{d=1}^{k} p_d(\sigma_d)
  = f(T).
\]
In practice, the target model is deterministic and the independence
assumption is approximate.  However, when the draft model is
well-calibrated, $f(T)$ serves as an effective proxy: maximizing $f$
consistently improves the realized acceptance length in our experiments
(Section~\ref{sec:experiments}).
\end{remark}

% ----------------------------------------------------------------------
\subsection{Monotonicity and Submodularity}
\label{sec:submodularity}
% ----------------------------------------------------------------------

\begin{proposition}[Monotone submodularity of $f$]
\label{prop:submodular}
The prefix-coverage objective $f \colon 2^{\mathcal{C}} \to \mathbb{R}_{\geq 0}$
is monotone non-decreasing and submodular.
\end{proposition}

\begin{proof}
We express $f$ as a non-negative weighted coverage function.  Let
$\mathcal{U} = \bigcup_{k=1}^{D} \mathcal{V}^k$ be the universe of
all possible prefixes up to depth~$D$.  Each element
$\sigma \in \mathcal{U}$ has weight $w(\sigma) \geq 0$.  Each leaf
$\ell \in \mathcal{C}$ \emph{covers} the set
$S(\ell) = \{\pi_1(\ell), \dots, \pi_{D_\ell}(\ell)\}$.  Then
\[
  f(T)
  = \sum_{\sigma \in \mathcal{U}} w(\sigma)
    \cdot \mathbf{1}\!\bigl[\sigma \in \textstyle\bigcup_{\ell \in T} S(\ell)\bigr].
\]
This is a non-negative weighted set-cover function, which is a standard
instance of a monotone submodular function~\citep{nemhauser1978}.

\textbf{Monotonicity.}\;
Adding a leaf $\ell$ to $T$ can only enlarge
$\bigcup_{\ell' \in T} S(\ell')$, so $f(T \cup \{\ell\}) \geq f(T)$.

\textbf{Submodularity.}\;
For $A \subseteq B \subseteq \mathcal{C}$ and any $\ell \notin B$,
the marginal gain is
\[
  \Delta(\ell \mid A)
  = \sum_{k=1}^{D_\ell}
    \mathbf{1}\!\bigl[\pi_k(\ell) \notin \mathrm{Prefixes}_k(A)\bigr]
    \cdot w\!\bigl(\pi_k(\ell)\bigr).
\]
Since $A \subseteq B$ implies
$\mathrm{Prefixes}_k(A) \subseteq \mathrm{Prefixes}_k(B)$,
every indicator that is~$1$ for~$B$ is also~$1$ for~$A$.
Hence $\Delta(\ell \mid A) \geq \Delta(\ell \mid B)$,
which is the definition of submodularity.
\end{proof}

% ----------------------------------------------------------------------
\subsection{Approximation Guarantee}
\label{sec:guarantee}
% ----------------------------------------------------------------------

\begin{corollary}[Greedy $(1 - 1/e)$-approximation]
\label{cor:greedy}
Let $T^* = \arg\max_{|T| \leq M} f(T)$ be the optimal leaf set of
size at most~$M$.  The greedy algorithm that iteratively selects
\[
  \ell_{i+1}
  = \arg\max_{\ell \in \mathcal{C} \setminus T_i}
    \Delta(\ell \mid T_i),
  \qquad
  T_{i+1} = T_i \cup \{\ell_{i+1}\},
\]
starting from $T_0 = \emptyset$, produces a set $T_M$ satisfying
\[
  f(T_M) \;\geq\; \Bigl(1 - \frac{1}{e}\Bigr) \cdot f(T^*).
\]
\end{corollary}

\begin{proof}
This is a direct application of the classical result of
\citet{nemhauser1978} for maximizing a monotone submodular function
under a cardinality constraint, which applies because $f$ is monotone
submodular (Proposition~\ref{prop:submodular}) and
$f(\emptyset) = 0$.
\end{proof}

% ----------------------------------------------------------------------
\subsection{Algorithm}
\label{sec:algorithm}
% ----------------------------------------------------------------------

Our prefix-aware tree builder operates in two phases.

\paragraph{Phase~1: Candidate generation.}
We generate a pool $\mathcal{C}$ of candidate leaves using best-first
search ordered by cumulative log-probability
$\sum_{d=1}^{|\ell|} \log p_d(\sigma_d)$.  At each step, the
highest-probability frontier node is expanded with up to $K$ children
(the top-$K$ tokens at the next depth).  Expansion continues until the
pool reaches a budget of $|\mathcal{C}| \approx 3M$ candidates.  This
phase is identical to the best-first tree builder of
Section~\ref{sec:bestfirst} but with a larger pool to provide diverse
candidates for the selection phase.

\paragraph{Phase~2: Greedy submodular selection.}
From the candidate pool $\mathcal{C}$, we greedily select $M$~leaves
to maximize $f(T)$.  At each iteration, we select the candidate with
the largest marginal gain $\Delta(\ell \mid T_i)$ and update the
covered prefix sets.

We implement this using the \emph{lazy greedy} acceleration
\citep{minoux1978}: a max-heap stores upper bounds on marginal gains.
When a candidate is popped, its marginal gain is recomputed; if it
remains the largest, it is accepted; otherwise it is re-inserted with
the updated value.  By submodularity, marginal gains can only decrease
as $T_i$ grows, so the stale heap values are valid upper bounds.  Lazy
greedy returns the \emph{identical} solution as standard greedy and
therefore inherits the $(1 - 1/e)$ guarantee of
Corollary~\ref{cor:greedy}, while typically requiring fewer
marginal-gain evaluations in practice.

\paragraph{Trie packing.}
The selected leaves are inserted into a trie.  Shared prefixes are
stored once, producing a packed token sequence, position-id sequence,
and parent-index array suitable for a single target-model forward pass
with tree attention.

% ----------------------------------------------------------------------
\subsection{Computational Cost}
\label{sec:cost}
% ----------------------------------------------------------------------

Phase~1 (best-first expansion) performs $O(|\mathcal{C}| \log |\mathcal{C}|)$
heap operations.  Phase~2 (greedy selection) performs at most $O(NM)$
marginal-gain evaluations in the worst case, where $N = |\mathcal{C}|$
and $M$ is the leaf budget.  Each evaluation scans at most $D$~prefix
depths, giving worst-case $O(NMD)$.  In practice, lazy greedy
substantially reduces the number of evaluations because most
candidates' marginal gains drop below the current best after the first
few selections.

The dominant cost in speculative decoding is the target-model forward
pass over the packed trie, which scales with the number of trie
nodes.  The tree construction overhead (Phase~1 + Phase~2) is
negligible by comparison: it operates on CPU over $O(N)$ candidate
sequences of length at most~$D$, with no GPU synchronization.

% ======================================================================
% References (add to your main .bib file)
% ======================================================================
% @article{nemhauser1978,
%   author  = {Nemhauser, George L. and Wolsey, Laurence A. and Fisher, Marshall L.},
%   title   = {An analysis of approximations for maximizing submodular set functions---{I}},
%   journal = {Mathematical Programming},
%   volume  = {14},
%   number  = {1},
%   pages   = {265--294},
%   year    = {1978},
% }
%
% @article{minoux1978,
%   author  = {Minoux, Michel},
%   title   = {Accelerated greedy algorithms for maximizing submodular set functions},
%   journal = {Optimization Techniques},
%   pages   = {234--243},
%   year    = {1978},
%   publisher = {Springer},
% }
